%!TEX root = ../main.tex
\section{Идея 3: Дистилляция Дейкстры}

\subsection{Зачем нам дистилляция и как мы это задумали}

Граф Дейкстры работал отлично, но ради спортивного интереса и ускорения я решил проверить: \textit{а можно ли обучить компактную нейросеть-студента, которая за один прямой проход ($O(1)$) будет предсказывать идеальную интенцию напрямую в $\pi^\ell$, действуя как виртуальная Дейкстра?}

Чисто механически сеть работает так же, как и бейзлайн: принимает $(s_t, B(g))$ и выдаёт одну интенцию. Но благодаря обучению на траекториях учителя она не пытается идти сквозь стены, а сразу выбирает правильный вектор обхода.

\subsection{4-компонентный дифференцируемый лосс}

Для обучения студента $T_\theta(s_t, B(g)) \to \hat{z}_t \in \mathbb{S}^{d-1}$ я собрал датасет пар $(s_t, B(g)) \to (z^*_t, a^*_t)$ от учителя Дейкстры и объединил 4 слагаемых в дифференцируемый лосс:
\begin{equation}
\label{eq:distillation_loss}
\begin{split}
\mathcal{L}(\theta) = {} & \underbrace{(1 - \cos(\hat{z}_t, z^*_t)) + 0.1 \|\hat{z}_t - z^*_t\|^2}_{\mathcal{L}_{\text{BC}}} \\
& + 0.5 \cdot \underbrace{\|\pi^\ell(s, \hat{z}) - a^*\|^2}_{\mathcal{L}_{\text{Action}}} \\
& - 0.02 \cdot \underbrace{F(s, \hat{z})^\top \hat{z}}_{\mathcal{L}_{\text{Reach}}} - 0.05 \cdot \underbrace{\hat{z}^\top B(g)}_{\mathcal{L}_{\text{Goal}}}
\end{split}
\end{equation}
где $\mathcal{L}_{\text{BC}}$ выравнивает вектор интенции с учителем на сфере; $\mathcal{L}_{\text{Action}}$ дифференцируется сквозь замороженную политику $\pi^\ell$ и штрафует за ошибки в крутящих моментах моторов робота; $\mathcal{L}_{\text{Reach}}$ максимизирует меру физической проходимости пространства FB; а $\mathcal{L}_{\text{Goal}}$ удерживает проекцию вектора на глобальную цель.

\begin{figure*}[t]
    \centering
    %!TEX root = ../main.tex
\begin{tikzpicture}[>=stealth, font=\small, scale=0.88, every node/.style={transform shape}]
    % Входные данные
    \node[draw=blue!80, fill=blue!10, rounded corners=4pt, minimum width=2.2cm, minimum height=0.9cm, align=center] (obs) at (0, 2.2) {Состояние $s_t$\\$\mathbb{R}^{29}$};
    \node[draw=orange!80, fill=orange!10, rounded corners=4pt, minimum width=2.2cm, minimum height=0.9cm, align=center] (goal) at (0, -0.2) {Цель $B(g)$\\$\mathbb{R}^{128}$};

    % Энкодеры
    \node[draw=blue!80, fill=blue!20, rounded corners=4pt, minimum width=2.4cm, minimum height=0.9cm, align=center] (enc_s) at (3.2, 2.2) {State Encoder\\MLP + SwiGLU};
    \node[draw=orange!80, fill=orange!20, rounded corners=4pt, minimum width=2.4cm, minimum height=0.9cm, align=center] (enc_g) at (3.2, -0.2) {Goal Encoder\\MLP + SwiGLU};

    % Блок внимания и гейта
    \node[draw=purple!80, fill=purple!15, rounded corners=6pt, minimum width=3.2cm, minimum height=1.6cm, align=center] (gate) at (7.0, 1.0) {\textbf{Gated Cross-Attention}\\$Q = E_s, \;\; K, V = E_g$\\$H = \text{Attn} \odot \sigma(W_g [E_s, E_g])$};

    % Проекция на сферу
    \node[draw=green!70!black, fill=green!15, rounded corners=4pt, minimum width=2.5cm, minimum height=0.9cm, align=center] (proj) at (10.8, 1.0) {\textbf{Проекция на $\mathbb{S}^{d-1}$}\\$\hat{z} = \sqrt{d} \cdot \frac{z}{\|z\|}$};

    % Выход
    \node[draw=red!70!black, fill=red!15, rounded corners=4pt, minimum width=2.2cm, minimum height=0.9cm, align=center] (out) at (13.8, 1.0) {\textbf{Интенция}\\$\hat{z}_{\text{cmd}} \to \pi^\ell$};

    \draw[->, thick] (obs) -- (enc_s);
    \draw[->, thick] (goal) -- (enc_g);
    \draw[->, thick] (enc_s) -- (gate);
    \draw[->, thick] (enc_g) -- (gate);
    \draw[->, thick] (gate) -- (proj);
    \draw[->, thick] (proj) -- (out);
\end{tikzpicture}

    \caption{Схема дистиллированного планировщика Distilled JAX Gated-Attention с дифференцируемым физическим лоссом.}
    \label{fig:distilled_jax}
    \vspace{-2mm}
\end{figure*}

\subsection{Какая архитектура показала себя лучше всего?}

Я поэкспериментировал с 4 архитектурами нейросетей (табл.~\ref{tab:arch_comparison}): обычным перцептроном MLP, остаточными блоками ResNet-ECA (с 1D-канальным вниманием), Dense-ECA и Gated Cross-Attention.

\begin{table}[H]
    \centering
    \footnotesize
    \setlength{\tabcolsep}{4pt}
    \caption{Сравнение архитектур дистилляции на Medium}
    \label{tab:arch_comparison}
    \begin{tabular}{lcccc}
        \toprule
        \textbf{Архитектура} & \textbf{Парам.} & \textbf{Cos-Sim} & \textbf{SR (\%)} & \textbf{Время} \\
        \midrule
        Standard MLP & 312K & 0.891 & 78.0\% & 0.12 мс \\
        Dense-ECA & 458K & 0.914 & 82.0\% & 0.16 мс \\
        ResNet-ECA & 524K & 0.908 & 80.0\% & 0.15 мс \\
        \textbf{Gated Cross-Attn} & \textbf{680K} & \textbf{0.942} & \textbf{86.0\%} & \textbf{0.18 мс} \\
        \bottomrule
    \end{tabular}
\end{table}

Лучше всех оказался \textbf{Gated Cross-Attention}: он показал на бенчмарке $84.3 \pm 4.8\%$ успешности (даже немного обогнав Дейкстру) при инференсе всего $0.55$~мс на шаг, что почти в два раза быстрее Дейкстры. Но стоит учесть, что алгоритм Дейкстры работает на CPU, а нейросеть — на GPU.

